\chapter{Scripts Python}
\label{app:scripts}

Este apêndice lista os scripts que produzem os resultados da dissertação e
reproduz os trechos centrais do código. Os trechos são copiados dos scripts,
sem os comentários, com as quebras de linha ajustadas e com linhas não
essenciais omitidas (indicadas por \texttt{\# ...}). A ordem de
execução e a estrutura de pastas estão no Apêndice~\ref{app:pastas}.

\section{Scripts por etapa}

O sufixo de cada script de análise (T1, T5, T10 etc.) é o nome da pasta em que
ele grava as saídas, dentro de \texttt{outputs/}.

\begin{description}
    \item[Coleta.] \texttt{code/update\_btc\_prices.py} baixa o preço de
          fechamento diário do Bitcoin (par BTCUSDT) da API pública da Binance,
          e \texttt{code/fix\_btc\_gap\_2023\_03.py} completa, com a mesma
          fonte, o intervalo de 1º a 31 de março de 2023.
          \texttt{code/download\_dvol\_deribit.py} baixa o DVOL da API pública
          da Deribit.
    \item[Construção.] \texttt{code/build\_vrp\_dataset.py} confere a
          continuidade diária das séries, calcula a volatilidade histórica nas
          duas janelas e as duas medidas do prêmio;
          \texttt{code/build\_targets.py} acrescenta os retornos futuros e
          forma a amostra completa; \texttt{code/build\_ml\_dataset\_T4.py}
          forma a amostra de modelagem, com o alvo e as 17 variáveis
          explicativas; \texttt{code/regimes\_T9.py} estima os cortes e
          classifica os regimes.
    \item[Utilitários.] \texttt{code/split\_utils.py} gera as origens de
          previsão, os treinos embargados e as dobras da validação cruzada;
          \texttt{code/hac\_utils.py} estima o MQO com erro-padrão de
          Newey--West; \texttt{code/avaliacao\_utils.py} calcula o $R^2$ fora
          da amostra e os testes de Clark--West e de Diebold--Mariano.
    \item[Capítulo~\ref{chap:dados}.] \texttt{code/build\_desc\_stats\_T1.py}
          (estatísticas descritivas e testes de raiz unitária),
          \texttt{code/plot\_cap3\_T1.py} (séries, histogramas e distribuição
          anual), \texttt{code/plot\_vrp\_vs\_return\_T12.py} (dispersão entre
          a proxy e o retorno de 30 dias), \texttt{scripts/test\_h1\_bvrp\_mean.py}
          (teste de H1) e \texttt{code/descritivas\_C3\_regimes.py} (prêmio por
          regime).
    \item[Capítulo~\ref{chap:previsao_bvrp}.] \texttt{code/previsao\_bvrp\_T5.py}
          (previsões dos cinco modelos e das referências, validação cruzada e
          placebo) e \texttt{code/interacao\_regime\_T9.py} (interações com o
          regime).
    \item[Capítulos~\ref{chap:retornos_linear}
          e~\ref{chap:retornos_nao_linear}.] \texttt{code/retorno\_bvrp\_T10.py}
          (modelo linear e bootstrap em dois níveis),
          \texttt{code/retorno\_bvrp\_T11.py} (modelo em árvore, testes
          conjuntos e quebras) e \texttt{code/descritivas\_C7.py} (retorno por
          regime e sinal do retorno).
    \item[Apêndice~\ref{app:estrategias}.] \texttt{code/estrategias\_A1.py}.
    \item[Geradores.] \texttt{code/publicar\_cap5.py},
          \texttt{code/publicar\_cap6.py}, \texttt{code/publicar\_cap7.py} e
          \texttt{code/publicar\_apendice.py} gravam as tabelas e as figuras
          dos Capítulos~\ref{chap:previsao_bvrp} a~\ref{chap:retornos_nao_linear}
          e do Apêndice~\ref{app:estrategias} a partir das saídas gravadas, sem
          reestimar modelos.
\end{description}

\section{Volatilidade histórica e prêmio}

Em \texttt{code/build\_vrp\_dataset.py}, a volatilidade histórica de 30 dias é
calculada nas duas janelas, em \% a.a., e as duas medidas do prêmio saem da
diferença com a IV:

\begin{verbatim}
df["ret"] = np.log(df["close"] / df["close"].shift(1))
window = 30
df["rv_30d"] = (
    (df["ret"] ** 2)
    .rolling(window=window)
    .mean()
    * 365
) ** 0.5 * 100
fwd = pd.api.indexers.FixedForwardWindowIndexer(window_size=window)
df["rv_30d_fut"] = (
    (df["ret"].shift(-1) ** 2)
    .rolling(window=fwd, min_periods=window)
    .mean()
    * 365
) ** 0.5 * 100
# ...
df = pd.merge(df_price, df_iv, on="date", how="inner")
# ...
df["vrp_30d"] = df["rv_30d"] - df["iv_30d"]
df["bvrp_30d_fut"] = df["rv_30d_fut"] - df["iv_30d"]
\end{verbatim}

\section{Origens de previsão e embargo}

Em \texttt{code/split\_utils.py}, cada origem $t$ estima o modelo com as
posições até $t - h$ e prevê as datas até a véspera da origem seguinte. Na
janela expansiva, o treino começa na primeira observação; na móvel, tem
tamanho fixo:

\begin{verbatim}
t0 = n_inicial - 1 + h0
# ...
origens = []
for t in range(t0, n, freq):
    fim = t - h
    inicio = 0 if janela == "expansiva" else fim - n_movel + 1
    origens.append(Origem(t=t,
                          treino=np.arange(inicio, fim + 1),
                          teste=np.arange(t, min(t + freq, n))))
\end{verbatim}

Com os valores usados, $n_{\text{inicial}} = 730$, $h_0 = 60$ e
$\textit{freq} = 30$: a primeira origem é 21/06/2023, e o período fora da
amostra tem 987 datas em 33 origens (Seção~\ref{sec:met-amostras}).

\section{Erro-padrão de Newey--West}

Todas as regressões usam a mesma função de \texttt{code/hac\_utils.py}, com
\texttt{maxlags} $= h + 1$, núcleo de Bartlett, sem correção de amostra
pequena e com inferência pela distribuição normal:

\begin{verbatim}
res = sm.OLS(y, exog).fit(
    cov_type="HAC",
    cov_kwds={"maxlags": L,
              "weights_func": weights_bartlett,
              "use_correction": False},
    use_t=False,
)
\end{verbatim}

\section{Avaliação fora da amostra}

Em \texttt{code/avaliacao\_utils.py}, o $R^2$ fora da amostra compara a soma
dos erros quadráticos do modelo com a da referência, e o teste de Clark--West
regride a diferença de perdas ajustada numa constante, com o mesmo erro-padrão
de Newey--West:

\begin{verbatim}
return float(1 - np.sum((y - prev) ** 2)
             / np.sum((y - prev_ref) ** 2))
# ...
# pm: previsão do modelo; pr: previsão da referência
y, pm, pr = _arr(y, prev_modelo, prev_ref)
f = (y - pr) ** 2 - ((y - pm) ** 2 - (pr - pm) ** 2)
res = mqo_newey_west(f, h=h)
t = float(res.tvalues.iloc[0])
\end{verbatim}

\section{Conferências e testes}

Antes de gravar uma tabela, cada gerador \texttt{publicar\_*} recalcula as
estatísticas que ela reporta a partir das previsões e das séries gravadas e
interrompe a execução se a diferença em relação às saídas da análise passar de
$10^{-9}$. Os testes em \texttt{code/test\_*.py} saem com erro se alguma
verificação falhar. Entre outras, eles conferem:
\begin{itemize}
    \item que a VH prospectiva de $t$ é a retrospectiva de $t+30$
          (\texttt{test\_T1\_alinhamento.py});
    \item que nenhuma variável explicativa muda quando os dados brutos são
          cortados em $t$ (\texttt{test\_T4\_vazamento.py});
    \item o embargo de todas as origens, com as datas reais
          (\texttt{test\_split\_utils.py});
    \item a convenção do erro-padrão de Newey--West
          (\texttt{test\_hac\_utils.py});
    \item o corte fixo dos regimes, recalculado dos dados brutos cortados no
          fim da primeira janela (\texttt{test\_regimes\_T9.py});
    \item as métricas fora da amostra, contra cálculo manual
          (\texttt{test\_T5.py});
    \item as regressões dos retornos e a ausência de informação futura nas
          estratégias (\texttt{test\_T10.py}, \texttt{test\_T11.py} e
          \texttt{test\_A1.py}).
\end{itemize}
